\documentclass[12pt]{article}
\usepackage[utf8]{inputenc}
\usepackage[T1]{fontenc}
\usepackage{hyperref}
\title{Zinza Language Resources for K-12 Curriculum Development}
\author{Compiled by Claude Sonnet 5.5 (claude-sonnet-5-5) for Andr\'as Kornai}
\date{}
\begin{document}
\maketitle

\section{Introduction}
Zinza (ISO 639-3: zin; Glottocode: zinz1238; Kizinza) is a Bantu language spoken on the southwest shores and islands of Lake Victoria in Tanzania (Mwanza region). Dialects named are Kula and Longo. Speaker estimates vary, from about 140,000 to 205,000, and Joshua Project reports 453,000 for the people group. No dictionary was found in search results.

\section{Primary Educational Resources}
No primers, textbooks or school programmes in Zinza were found in search results. A 2000 booklet of 105 Zinza proverbs with Swahili and English translations may be useful as reading material.
\begin{itemize}
\item African Proverbs, Sayings and Stories (Tangaza University): \url{https://afriprov.tangaza.ac.ke/?p=464}
\end{itemize}

\section{Linguistic and Reference Materials}
\begin{itemize}
\item Numeral-system data sheet for Zinza (Mwanza Region): \url{https://lingweb.eva.mpg.de/channumerals/Zinza.htm}
\item Joshua Project, language page for zin: \url{https://www.joshuaproject.net/index.php/languages/zin}
\end{itemize}

\section{Religious and Community Resources}
Scripture Earth lists a Zinza New Testament, an app, audio and the JESUS Film. Joshua Project notes older scripture portions from 1930 in a hard-to-read alphabet, since revised. A Wycliffe project profile describes the translation work.
\begin{itemize}
\item Scripture Earth, Zinza: \url{https://www.scriptureearth.org/00eng.php?iso=zin}
\item Wycliffe, Zinza 2 project profile: \url{https://www.wycliffe.org/Main%20Uploads/Projects/PDF/Zinza%202%20Profile%20AMENDED%20WEB.pdf}
\item Joshua Project, Zinza in Tanzania: \url{https://joshuaproject.net/people_groups/16108}
\end{itemize}

\section{Digital and Audio Resources}
\begin{itemize}
\item Global Recordings Network, Zinza (audio Bible stories, MP3 and slideshow packs): \url{https://globalrecordings.net/en/language/zin}
\end{itemize}

\section{Curriculum Development Context}
Tanzanian schools use Swahili and English. Zinza has an established orthography, as shown by the New Testament, and a translation community, but no known educational materials. A curriculum could begin with the scripture texts, the proverb collection and audio recordings, and would need new reading materials and a dictionary.

\end{document}